% 2009 年山东卷理科第 4 题：三视图（圆柱与其上四棱锥）。
% 来源 PDF：普通高考/2009/2009山东理.pdf，SHA256 6c3546f276b37d6ff4fa1a8cdf77e8356b8a4c0e713e9e5485bb47fcd5ed660f
% 原页：第 1 页（A3 三栏）第 1 栏中部，印刷页 545。
% 重绘依据：正视图为 2×2 矩形上方接顶角三角形（斜棱标 2），侧视图同形；
%   俯视图为半径 1 的圆内接正方形（对角线为两条互相垂直的直径）；尺寸按原图标注。
\begin{tikzpicture}[scale=0.8, semithick, every node/.style={inner sep=1.3pt, font=\small}]
  % 正视图
  \begin{scope}[shift={(0,0)}]
    \draw (0,0) rectangle (2,2);
    \draw (0,2) -- (1,3.0) -- (2,2);
    \node[anchor=east] at (-0.12,1) {$2$};
    \node[anchor=north] at (1,-0.12) {$2$};
    \node[anchor=south east] at (0.42,2.55) {$2$};
    \node[anchor=north] at (1,-0.42) {正（主）视图};
  \end{scope}
  % 侧视图
  \begin{scope}[shift={(4.2,0)}]
    \draw (0,0) rectangle (2,2);
    \draw (0,2) -- (1,3.0) -- (2,2);
    \node[anchor=north] at (1,-0.12) {$2$};
    \node[anchor=south east] at (0.42,2.55) {$2$};
    \node[anchor=north] at (1,-0.42) {侧（左）视图};
  \end{scope}
  % 俯视图
  \begin{scope}[shift={(1.0,-4.6)}]
    \draw (1,1) circle (1);
    \draw (0,1) -- (2,1) -- (1,2) -- (1,0) -- cycle;
    \node[anchor=north] at (1,-0.18) {俯视图};
  \end{scope}
\end{tikzpicture}
